\section{Artifacts, verification status and remaining problem}\label{sec:release}

\subsection{Three distinct roles of computation}

The artifact separates complete proof kernels, explicit graph witnesses,
and bounded implementation tests. The first role is logically part of
the computer-assisted proofs. A1 uses the complete safe-core relation
tables, 163 sparse-core representatives, and a 4,256-case weak two-mark
kernel, with fixed-point and path-saturation arguments for unbounded
length. The clique interface uses 836 complete small-order offset bases
and its two-vertex induction. The decorated-path interface uses all
93 endpoint comparison types, all 96 transition rows, and the displayed
lower-bound and saturation identities. These domains are finite because
their equivalence reductions are proved, not because experiments stopped
at a chosen graph order.

The second role is independently checking submitted weight assignments.
The release verifier rebuilds red and blue degrees from the actual
edge-copy counts, including physical offsets when present, verifies all
required comparisons, checks any recorded final degrees, and confirms
preservation of supplied boundary weights. Large stored region witnesses
can therefore be audited without repeating their discovery search.
The 13-regular descent barrier is supported by a fixed integer table
whose row sums and neighbour counts are checked directly; its arbitrary
fibre sizes are justified by the explicit matching construction.

The third role is testing the constructive implementations on graphs,
offsets and lengths beyond the finite kernels. These checks include
small graph atlases, random compressed graphs, long paths and return
paths, prescribed-root cliques, and several complete dense-region
examples. They audit code, rather than replace any infinite proof.
Original research-run reports are retained separately from reports of
the new release run. The provenance ledger identifies frozen original
notes and the script dependencies. Optional factor searches use
NetworkX and sometimes Z3. Their timeouts are reported as unfinished
searches; a timeout is never treated as an infeasibility certificate.

\subsection{Reproducing the release}

The default complete-certificate command, from an extracted archive, is
\begin{center}\code{python -I -B verify.py}\end{center}
It uses the Python standard library. It invokes the A1 certificate
verifier and its independent local recalculation, then the additional
clique and decorated certificates, fixed integer tables, stored graph
witnesses and independent arithmetic checks. It also runs all four
portable E1--E4 audits without importing their constructors or solvers.
Those audits recount 972 actual copies in the saved whole-graph witnesses
and rederive the submitted rows, local relations, parity cuts and
strategy obstructions. Do not run Python with
optimisation enabled: the certificate checkers deliberately use
assertions. A JSON report can be written using the option
\texttt{-{}-output validation/certificates.json}.

For constructive implementation tests, install the pinned dependencies
in \code{requirements.txt} and run
\begin{center}\code{python -B run_checks.py}\end{center}
This runs the released construction checkers and saves a bounded report.
The A1 command-line constructor remains available as
\code{construct.py}, with examples and an independent copy-degree
checker in the root directory. Other constructors and their scope
checks are in \code{artifact/}; the coverage ledger points to each
public implementation. A method-not-applicable response is not a
counterexample to the mathematical conjecture.

The editable English source is \code{manuscript/main.tex} with included
section files and an explicit bibliography. It can be compiled by a
standard LaTeX installation or Tectonic. No user-specific absolute path
is required by the released proof data or verifier. The manifest records
SHA-256 digests of the released files. The final validation report records
fresh extraction, certificate execution, constructor checks, PDF rebuild
and source-reference diagnostics. Checks of the archive are local;
no repository deposit, Zenodo upload, DOI allocation, or journal
submission is implied.

\subsection{Use of AI and research status}

This manuscript and its research artifacts were developed with extensive
assistance from OpenAI Codex, including mathematical exploration, proof
drafting, programming, calculations, internal review and English editing.
Automated checks and internal reviews are documented as such; they
are not independent expert peer review. No Lean or other proof-assistant
verification is claimed.

The artifact attributes the pre-existing graph-family results and the
classical degree-list, factor, matching and interpolation tools used in
the proofs. Its grouping of results records this research programme's
contributions and boundaries; it is not an assertion that every grouped
lemma first appears here. A literature-screening record is supplied as
provenance, with no guarantee of exhaustive historical priority.
The manuscript, documentation and finite certificate data are licensed
under Creative Commons Attribution 4.0 International (CC BY 4.0).
The accompanying Python source code is licensed under the MIT licence.
The licence-scope file and full licence texts accompany the release.

The remaining mathematical problem is the original arbitrary-graph
conjecture. None of the positive results forces a general graph to have
a sparse, homogeneous, clique, parity-separated, high-degree or dense
region structure. Target optimisation is exact only after a target is
fixed, and the algebraic route lacks a general nonvanishing theorem.
The E1--E4 partition conditions and mixed-port certificates have not been
proved to exist for every graph. In particular, a parity-subgraph
solution can saturate a forbidden local component, and root-only joins
can exclude valid colourings using even components.
Minimal-counterexample restrictions still permit arbitrary large,
nonregular, minimum-degree-three cores. The explicit obstructions show
why several universal parity and monotone-repair assertions cannot
fill that gap. A1 covers the full subcubic range; the later conditions
extend beyond that range but do not combine into a proof for every graph.
